% Corrected French for printed p.161.
% Corrections: restore O* as the resolved sheaf; rationnelle.
\SourcePage{161}

\LayerTarget{tohoku:construction:divisor-sequence:p2}
\noindent $\mathbf D$ est flasque, car une section de $\mathbf D$ sur un ouvert non
vide $U\subset X$ est une combinaison linéaire formelle
$\sum n_iV_i$ d'hypersurfaces irréductibles $V_i$ dans $U$, et est donc la
restriction de la section $\sum n_iV_i$ de $\mathbf D$ sur $X$. Par suite,

\LayerTarget{tohoku:eq:divisor-resolution}
\[
0\longrightarrow\mathbf O^*\longrightarrow\mathbf K^*\longrightarrow
\mathbf D\longrightarrow0\longrightarrow0\longrightarrow\cdots
\]
est une résolution de $\mathbf O^*$ par des faisceaux flasques
($\mathbf K^*$ est flasque puisqu'il est constant et que $X$ est
irréductible, cf. fin de 3.3). On en conclut les valeurs des
$H^p(X,\mathbf O^*)$ (on omet le signe $\Phi$ quand on prend pour $\Phi$ la
famille de toutes les parties fermées de $X$) :

\LayerTarget{tohoku:eq:divisor-class-cohomology}
\[
H^1(X,\mathbf O^*)=
\Gamma(\mathbf D)/\operatorname{Im}\bigl(\Gamma(\mathbf K^*)\bigr)
\]
$={}$ groupe des classes de diviseurs modulo les diviseurs principaux (fait
bien connu, et qui résulte d'ailleurs aussitôt de la suite exacte de
cohomologie), et $H^i(X,\mathbf O^*)=0$ si $i\geq2$, (résultat que j'avais
obtenu d'abord, de façon beaucoup moins simple, par la méthode des
recouvrements de Čech). On notera d'ailleurs que ce résultat s'étend sans
changement au cas où $X$ est un \og schéma de variété\fg{}, au sens de
[5 bis], plus généralement au cas de \og variétés arithmétiques\fg{},
définies par \og recollement\fg{} à partir de \og spectres\fg{} d'anneaux
commutatifs [8]. L'application ci-dessous peut aussi se formuler dans le
cadre des variétés arithmétiques :

\LayerTarget{tohoku:prop:3.4.1}
\textsc{Proposition 3.4.1.} --- \textit{Soit $X$ une variété algébrique
irréductible (sur un corps algébriquement clos $k$) dont les anneaux locaux
sont factoriels (par exemple une variété sans singularités), alors on a
$H^i(X,\mathbf O^*)=0$ pour $i\geq2$. Si $E$ est un fibré algébrique
localement trivial sur $X$, de groupe structural le groupe projectif
$\operatorname{GP}(n-1,k)$ (cf. [20]), alors $E$ est isomorphe à l'espace
fibré associé à un espace fibré algébrique localement trivial de groupe
structural le groupe linéaire $\operatorname{Gl}(n,k)$.}

\LayerTarget{tohoku:proof:3.4.1}
Pour tout groupe algébrique $G$, désignons par $\mathbf O(G)$ le faisceau de
groupes des germes d'applications régulières de $X$ dans $G$. Alors la
première assertion de la proposition à démontrer s'écrit
$H^i(X,\mathbf O(k^*))=0$ pour $i\geq2$, et a déjà été prouvée, la deuxième
s'écrit, en utilisant les notions et la terminologie développée dans [11] :
l'application canonique

\LayerTarget{tohoku:eq:pgl-lift-map}
\[
H^1\bigl(X,\mathbf O(\operatorname{Gl}(n,k))\bigr)
\longrightarrow
H^1\bigl(X,\mathbf O(\operatorname{GP}(n-1,k))\bigr)
\]
est surjective. Pour le prouver, considérons la suite exacte de groupes
algébriques

\LayerTarget{tohoku:eq:central-extension}
\[
e\longrightarrow k^*\longrightarrow\operatorname{Gl}(n,k)
\longrightarrow\operatorname{GP}(n-1,k)\longrightarrow e,
\]
où le premier homomorphisme est l'isomorphisme naturel de $k^*$ sur le centre
de $\operatorname{Gl}(n,k)$. On constate facilement que la fibration de
$\operatorname{Gl}(n,k)$ par le sous-groupe $k^*$ est localement triviale
(i.e. il existe une section rationnelle), donc la suite exacte précédente
donne naissance à une suite exacte de faisceaux

\LayerTarget{tohoku:eq:sheaf-central-extension}
\[
e\longrightarrow\mathbf O(k^*)
\longrightarrow\mathbf O(\operatorname{Gl}(n,k))
\longrightarrow\mathbf O(\operatorname{GP}(n-1,k))
\longrightarrow e,
\]
où $\mathbf O(k^*)$ est dans le centre de
$\mathbf O(\operatorname{Gl}(n,k))$. La proposition résulte alors de
$H^2(X,\mathbf O(k^*))=0$, et du corollaire au résultat suivant, qui
généralise [11, prop. 5.7.2, corollaire], où nous étions obligés de faire des
hypothèses de paracompacité :

\LayerTarget{tohoku:prop:3.4.2:p1}
\textsc{Proposition 3.4.2.} --- \textit{Soient $X$ un espace topologique,
$e\longrightarrow F\longrightarrow G\longrightarrow H\longrightarrow e$ une}

\TohokuContinuation{162}
